\subsubsection{Como o algoritmo funciona}

\indent \qquad O Insertion Sort, ou ordenação por inserção, é um dos métodos de
ordenação mais simples que existem. Ele funciona como quem organiza cartas na
mão: as cartas da esquerda já estão em ordem e, a cada carta nova, procura-se o
lugar dela nesse trecho, afastando as demais para abrir espaço. É com essa mesma
analogia que Cormen \textit{et al.} (2012) apresentam o método.

\qquad O algoritmo percorre o vetor da segunda posição até a última. A cada
passagem vale uma regra: tudo o que está à esquerda da posição atual já se
encontra ordenado. Cormen \textit{et al.} (2012) chamam essa regra de invariante
de laço, e é dela que sai a prova de que o algoritmo funciona.

\qquad Uma passagem faz o seguinte. O elemento da posição atual é copiado para
uma variável auxiliar, o que deixa a célula livre. Esse valor é então comparado
com os elementos à esquerda, um a um, da direita para a esquerda. Todo elemento
maior que ele é empurrado uma casa para a direita. A comparação para assim que
aparece um elemento menor ou igual, ou quando o início do vetor é alcançado. O
valor guardado no auxiliar entra na casa que sobrou. Ao fim da passagem, o
trecho ordenado ganhou um elemento; ao fim de todas, ele é o vetor inteiro.

\qquad Duas características do método importam para a análise que vem adiante. A
primeira é que ele ordena \textit{in-place}: não usa um segundo vetor. Fora o
auxiliar e os índices, nada cresce junto com o tamanho da entrada, e por isso o
espaço extra é $O(1)$. A segunda é que ele é estável: um elemento só é empurrado
quando é estritamente maior que o valor a ser inserido, de modo que dois
elementos iguais nunca trocam de ordem entre si (ZIVIANI, 2011).

\subsubsection{Exemplo passo a passo}

\indent \qquad A Tabela \ref{tab:exemplo_insertion} mostra esse processo em um vetor de
cinco elementos, exibindo o estado dele ao fim de cada passagem. Nela, a
passagem $j = 3$ não move nada, enquanto a passagem $j = 4$ move três elementos.
O trabalho de cada passagem não é fixo: depende de quão longe o elemento está do
seu lugar definitivo. É dessa variação que nascem o melhor e o pior caso.

\begin{spacing}{1.0}
	\begin{table}[htbp]
		\centering
		\caption{Insertion Sort aplicado ao vetor [7, 3, 9, 2, 5]. Em negrito, o trecho já ordenado ao fim de cada passagem.}
		\label{tab:exemplo_insertion}
		\begin{tabular}{c|ccccc|l}
			\hline
			\textbf{Passagem} & \multicolumn{5}{c|}{\textbf{Vetor}} & \textbf{O que ocorreu} \\
			\hline
			inicial & \textbf{7} & 3 & 9 & 2 & 5 & apenas o 7 está ordenado \\
			$j = 2$ & \textbf{3} & \textbf{7} & 9 & 2 & 5 & o 7 desloca; o 3 vai ao início \\
			$j = 3$ & \textbf{3} & \textbf{7} & \textbf{9} & 2 & 5 & o 9 já é maior que o 7 \\
			$j = 4$ & \textbf{2} & \textbf{3} & \textbf{7} & \textbf{9} & 5 & o 2 percorre três elementos \\
			$j = 5$ & \textbf{2} & \textbf{3} & \textbf{5} & \textbf{7} & \textbf{9} & o 5 para logo após o 3 \\
			\hline
		\end{tabular}
	\end{table}
\end{spacing}

\subsubsection{Pseudocódigo e custo por linha}

\indent \qquad A Tabela \ref{tab:custo_insertion} apresenta o algoritmo em
pseudocódigo e associa a cada linha o seu custo. O termo $c_k$ é o custo de
executar uma vez a linha $k$, e $t_j$, o número de vezes que o teste do laço
interno é avaliado na passagem $j$. Os dois reaparecem na Seção
\ref{sec:analise_insertion}.

\begin{spacing}{1.0}
	\begin{table}[H]
		\centering
		\caption{Pseudocódigo do Insertion Sort, com o custo e o número de execuções de cada linha.}
		% [H] em vez de flutuante: como flutuante o LaTeX a empurrava para depois da quebra
		\label{tab:custo_insertion}
		\renewcommand{\arraystretch}{1.3}
		\begin{tabular}{rlcc}
			\toprule
			\multicolumn{4}{c}{\textbf{\textit{INSERTION-SORT(A)}}} \\
			\midrule
			& \textbf{PSEUDOCÓDIGO} & \textbf{CUSTO} & \textbf{VEZES} \\
			\midrule
			1- & \texttt{for j $\leftarrow$ 2 to length(A) do} & $c_1$ & $n$ \\
			2- & \texttt{\hspace*{1em}key $\leftarrow$ A[j]} & $c_2$ & $n-1$ \\
			3- & \texttt{\hspace*{1em}i $\leftarrow$ j - 1} & $c_3$ & $n-1$ \\
			4- & \texttt{\hspace*{1em}while i $>$ 0 and A[i] $>$ key do} & $c_4$ & $\sum\limits_{j=2}^{n} t_j$ \\
			5- & \texttt{\hspace*{2em}A[i + 1] $\leftarrow$ A[i]} & $c_5$ & $\sum\limits_{j=2}^{n}(t_j - 1)$ \\
			6- & \texttt{\hspace*{2em}i $\leftarrow$ i - 1} & $c_6$ & $\sum\limits_{j=2}^{n}(t_j - 1)$ \\
			7- & \texttt{\hspace*{1em}A[i + 1] $\leftarrow$ key} & $c_7$ & $n-1$ \\
			\bottomrule
		\end{tabular}
	\end{table}
\end{spacing}

\subsubsection{Implementação em linguagem C}

\indent \qquad A implementação usada neste trabalho é a versão tradicional apresentada
em sala de aula, escrita em linguagem C e reproduzida no Código
\ref{cod:insertion_sort}.

\begin{spacing}{1.0}
\begin{lstlisting}[caption={Implementação do Insertion Sort em linguagem C.},label={cod:insertion_sort}]
	
    int i, j, aux;

    for (i = 1; i < tamanho; i++) {

        aux = vetor[i];
        j = i - 1;

        while (j >= 0 && vetor[j] > aux) {
            vetor[j + 1] = vetor[j];
            j--;
        }

        vetor[j + 1] = aux;
    }

\end{lstlisting}
\end{spacing}
